\scanpage{4}
% archive foliation: 3
11) Caractérisation des \deleted{\unsure{équiv}}\unsure{SQ|W}-fibrations dans $Cat$

i.e. $X'\longrightarrow X$ tel que le chgt. de base transforme équiv. faibles en équiv. faibles --- il suffit que ce soit OK pour les inclusions $e_Y\longrightarrow Y$, $e_Y$ objet final de $Y$. Est-il suffisant que l'on le suppose pour \inserted{tout} $\Delta_0\longrightarrow\Delta_n\longrightarrow X$ ?

\margin{\diagram{\begin{tikzcd}[row sep=small,column sep=small]
Z' \arrow[r,dash] \arrow[d] & Z \arrow[d] \\
Y' \arrow[r,dash] \arrow[d] & Y \arrow[d] \\
X' \arrow[r,dash] & X
\end{tikzcd}}}

12) Introduction de la condition
\diagram{\begin{tikzcd}[row sep=small,column sep=small]
X' \arrow[r,"g'"] \arrow[d,"f'"'] & X \arrow[d,"f"] \\
Y' \arrow[r,"g"'] \arrow[ur,phantom,"\mathrm{cart}"] & Y
\end{tikzcd}}

$f$ fib, $g\in\underline W\implies g'\in\underline W$ dans axiomatique de Quillen.

13) Is a \deleted{tot} 0-connected topos also totally $\infty$-connected? \inserted{and conversely} Is it enough that product of two \inserted{0-}connected objects is 0-connected (or of \inserted{2} $\underline W$-aspheric objects is $\underline W$-aspheric)?

14) \uline{\unsure{Homs}} in $(Hot)$ and diagram variants\ldots

% Reused typed sheet, with handwritten mathematical symbols; orientation 180 degrees.
\deleted{4 janvier 1973}

\deleted{Nikias Stavroulakis\\
99, rue de Sèvres\\
92100 Boulogne}

\deleted{avec $2\mu=3\times10^5\,\mathrm{cm}$. La valeur de la constante $\rho_0$ n'est pas connue, mais les phénomènes connus conduisent à accepter une valeur de $\rho_0$ d'ordre de grandeur petit relativement aux distances qui entrent en jeu. En effet, si l'on tient compte de ce que $2\mu\log\left(\sqrt{\frac{\unsure{g}}{2\mu}}+\sqrt{\frac{\unsure{g}}{2\mu}-1}\right)$ est, pour ces distances, très petit relativement à $\sqrt{\unsure{g}(\unsure{g}-2\mu)}$, on aura approximativement $\unsure{g}(\rho)=\rho$ lorsque $\rho_0$ a une petite valeur. Nous pouvons donc utiliser (6.3) avec $y=\frac1\rho$ et chercher, comme d'habitude, les solutions approximatives de l'équation qui s'obtient en dérivant par rapport à $\varphi$, à savoir}
\[
\deleted{\frac{d^2y}{d\varphi^2}+y=\frac{\mu}{E_2^2}+3\mu y^2\,,}
\]
\deleted{ce qui donne les valeurs numériques connues concernant la déviation des rayons lumineux et le déplacement du périhélie de Mercure. Ces valeurs ne dépendent donc pas de la valeur exacte de $\rho_0$, mais seulement de l'ordre de grandeur de $\rho_0$.}

15) Topos relt. \unsure{clg}, défini par un \uline{champ} sur $B$.

% Printed pagination of the reused sheet, not an author page number of these notes.
-- 18 --
